\reservewrapspace{28\baselineskip}
\section{System Overview}
\label{sec:chains}

\subsection{Derivation Pipeline}

\begin{wrapfigure}{r}{0.48\textwidth}
\centering
% Shared vector source for the paper and standalone diagram.
\begin{tikzpicture}[node distance=3mm,
  process/.style={stage,text width=4.5cm,font=\sffamily\small,minimum height=7mm}]
\node[process,draw=origblue] (extract) {Extract dependent reduction\\chains (\S\ref{sec:input-chains})};
\node[process,below=of extract] (judge) {Establish a single-pass form\\State-count upper bound (\S\ref{sec:existence})};
\node[process,below=of judge,draw=auxpurple] (candidate) {Reference tuple from source states\\Synthesise auxiliary states\\(\S\ref{sec:selection}--\ref{sec:synthesis})};
\node[process,below=of candidate] (derive) {Construct join through recovery\\Verify correction formulas\\(\S\ref{sec:recovery}--\ref{sec:rebase})};
\node[process,below=of derive] (check) {Emit and check fused program\\(\S\ref{sec:accuracy}--\ref{sec:impl})};
\node[process,below=of check,draw=fixgreen] (output) {Accept single-pass rewrite};
\foreach \a/\b in {extract/judge,judge/candidate,candidate/derive,derive/check,check/output}
  \draw[flow] (\a) -- (\b);
\draw[flow,draw=auxpurple] (derive.west) -- ++(-0.35,0) |- (candidate.west);
\node[lbl,rotate=90,anchor=south] at ([xshift=-3.5mm]derive.north west) {retry};
\node[lbl,below=2mm of output] {Unaccepted chains remain unchanged};
\end{tikzpicture}

\caption{End-to-end derivation.}
\label{fig:pipeline}
\end{wrapfigure}
\sys 以 TVM Tensor Expressions（TE）表示的程序及其声明的输入域为输入。
如 \cref{fig:pipeline} 所示，它首先提取沿同一轴执行的依赖归约链（\cref{sec:input-chains}），再为各条链推导单遍替代程序。
对于每条提取出的链，代数分析将每个归约的上下文与依赖输入的项分离。
得到的参考状态可以独立合并并确定输出，由此建立状态数上界，作为状态选择的默认限制（\cref{sec:existence}）。

执行其展开公式可能丢失源程序的中心化和缩放，因此选择过程以原有归约结果作为维护状态的起点。
若某个源结果在合法非空分区上无定义，则尝试提取导致定义域问题的上下文、保留其系数，以获得有定义的表示。
为发现缺失的信息，分析器推导将局部上下文替换为更大分区的上下文时，每个保留结果如何变化。
这些修正的系数提供候选辅助状态。
过程首先检查源状态自身，然后检查加入全部可用修正系数后的元组，即重定基闭包。
只有这两个元组都无法在状态数限制内成功时，语法引导合成才搜索这些特征的组合，随后将参考状态也加入搜索。
因此，修正分析可以直接提供所需辅助状态；当这些初始选择不足以满足要求时，合成进一步寻找紧凑的组合（\cref{sec:selection,sec:synthesis}）。

对于每个提出的维护元组，验证过程构造表达式，恢复用于表示其状态及所需输出的完整参考元组，并证明其恒等关系及在合法分区上的有定义性。
由此定义符号形式的合并操作：恢复两个分区的参考状态，分别组合，再将结果表示为维护状态。
验证确立合并后的元组等于并集的状态；证明未成功时，在搜索限制内返回状态选择。
随后，将已验证的合并重写为维护状态上的直接修正，尽可能保留中心化和缩放，以保持源程序的数值行为（\crefrange{sec:recovery}{sec:rebase}）。

后端生成 TE 程序，为每个输入元素计算叶状态元组，合并这些元组，再通过收尾计算产生所需输出。
针对归约器的静态检查与重建程序上的差分数值测量决定是否接受这一浮点实现（\cref{sec:accuracy,sec:impl}）。
若未找到通过验证的构造，或代码生成、数值检查失败，则保留源归约链。


\subsection{Input Reduction Chains}
\label{sec:input-chains}

提取出的链包含沿同一轴执行的归约及其间计算。
\Cref{fig:grammar} 给出了归约体的语言；以下定义说明链内依赖及其在分区上的语义。

\begin{figure}[htbp]
\small
\begin{align*}
f \;::=\;& x \mid \ctx \mid c \mid N
 \mid f + f \mid f - f \mid f \times f \mid f \div f
 \mid \exp f \mid \log f \mid \sqrt{f} \mid f^{k} \mid |f| \mid \max(f,f) \mid \min(f,f)\\
m \;::=\;& \textstyle\bigoplus_{j} f \mid g(\ctx_1,\dots,\ctx_p)
\end{align*}
\caption{Reduction bodies $f$ and chain members $m$. $x$ ranges over the values a chain reads at position $j$ of the axis, $\ctx$ over earlier members, $c\in\R$ over constants, $k\in\mathbb{Z}$, $N$ is the extent of the axis, and $\oplus\in\{\Sigma,\max,\min\}$.}
\label{fig:grammar}
\end{figure}

\begin{definition}[Reduction chain]
\label{def:chain}
固定一个轴 $\Ax=\{0,\dots,N-1\}$；对每个位置 $j\in\Ax$，以元组 $x_j$ 表示归约链在位置 $j$ 从链外张量读取的值。
\emph{归约链}是按生产者先于使用者的顺序排列的成员序列 $m_1,\dots,m_n$。
成员可以是\emph{归约环节} $m=\bigoplus_{j\in\Ax} f_m(x_j,\ctx_m)$，其中 $\oplus\in\{\Sigma,\max,\min\}$；也可以是一个\emph{值} $m=g_m(\ctx_m)$。
\emph{上下文} $\ctx_m$ 是 $m$ 使用不含 $j$ 的索引读取的较早成员，即这些成员的最终值。
只要某个归约环节读取了上下文，就称该链\emph{具有依赖关系}。
\end{definition}

在贯穿全文的方差示例中，最终求和读取均值作为上下文。
\sys 从 TE 程序中提取归约链：TE 归约成为归约环节，链沿 $j$ 读取的逐元素张量则内联到归约体中。
其他轴上的归约，例如注意力中的 $QK^\top$，保留在链外并提供 $x_j$ 中的值。

为了讨论轴的一部分，我们只在该部分上对归约链求值。

\begin{definition}[Partial semantics]
\label{def:partial}
对于非空的 $S\subseteq\Ax$，当每个归约环节仅在 $S$ 上归约、每个上下文取其自身在 $S$ 上的值、作为数值读取的轴长度 $N$ 取 $|S|$ 时，成员 $m$ 的值记为 $m(S)$。
原程序计算的是 $m(\Ax)$。
\end{definition}

声明的输入域约束元素值。
合法分区是满足这些约束的输入的任意非空部分。
即使完整输入的输出有定义，某个原始成员仍可能在这样的分区上无定义；维护状态则必须在每个合法分区上都有定义。
单遍算法必须从逐个分区收集的信息中，在源输出有定义的地方得到 $m(\Ax)$。

\subsection{Single-Pass Reducer Contract}

生成的程序必须从每个元素初始化状态元组，在不重新访问输入的情况下合并这些元组，并恢复所需输出。
以下约定给出每次合并在实数语义下必须保持的关系。

\begin{definition}[Single-pass reducer]
\label{def:reducer}
非空分区上的函数元组 $\st=(\st_1,\dots,\st_k)$ 称为\emph{状态映射}。
固定非归约索引后，每个函数的值均为标量，给定归约器的元组大小 $k$ 固定，不随归约长度 $N$ 增长。
归约链的\emph{单遍归约器}是由叶状态 $\leaf(x_j)=\st(\{j\})$、二元操作 $\Mrg$ 和收尾计算 $\epi$ 构成的元组 $(\leaf,\Mrg,\epi)$，对任意不相交的非空分区 $\pa$ 和 $\pb$，满足
\begin{equation}
\label{eq:invariant}
\Mrg(\st(\pa),\st(\pb))=\st(\pa\dunion\pb),\qquad
\epi(\st(\Ax))=\bigl(m(\Ax)\bigr)_{m\text{ required}}.
\end{equation}
\emph{所需}成员是程序其余部分读取的成员。
\end{definition}

生成的归约从第一个叶状态开始，随后折叠其余叶状态。
没有数值状态代表空分区：执行时跳过空分区，且输入轴必须非空。
\cref{eq:invariant} 中的合并不变量将维护状态元组与原程序联系起来。

\begin{lemma}[Laws from the invariant]
\label{lem:laws}
如果不变量成立，则在两两不相交的非空分区状态上，$\Mrg$ 满足交换律和结合律。
对 $\Ax$ 的任意非空分区划分，以及这些分区上的任意二叉合并树，合并叶状态均得到 $\st(\Ax)$。
从第一个叶状态初始化的串行折叠也具有相同性质。
\end{lemma}

\begin{proof}
$\Mrg(\st(\pa),\st(\pb))=\st(\pa\dunion\pb)=\st(\pb\dunion\pa)=\Mrg(\st(\pb),\st(\pa))$，结合律同样由 $\dunion$ 的结合律推出。
其余结论分别通过对合并树和折叠长度的归纳得到。
\end{proof}

因此，每个非空分区可以独立计算自己的状态元组。
合并树组合这些元组，无需再次遍历各个分区。
这一并行结构来自不变量，而不依赖某个特定调度。

用程序演算的术语来说，$S\mapsto\st(S)$ 是从以 $\dunion$ 为运算的有限非空多重集到以 $\Mrg$ 为运算的状态的同态，即合并与顺序无关的列表同态~\citep{bird87,gibbons96}。
用数据库的术语来说，具有无上下文参考构造的链是一个\emph{代数型}聚合，可由固定数量的\emph{分配型}聚合计算得到~\citep{datacube}。
空分区由控制流处理，不属于代数合并的约定范围。

\paragraph{Problem.}
给定一个依赖归约链，建立单遍归约器的充分条件，构造归约器，证明 \cref{eq:invariant}，并且仅当浮点计算在检查的输入上不引入额外的非有限结果、测得误差相对于原程序保持在规定界限内时接受融合。
